\documentclass[a4paper,12pt]{article}

\usepackage[T1]{fontenc}
\usepackage[italian]{babel}
\usepackage{graphicx}
\usepackage{geometry}
\usepackage{tabularx}
\usepackage[table]{xcolor}
\usepackage[most]{tcolorbox}
\usepackage{fancyhdr}
\usepackage[hidelinks]{hyperref}

\newcommand{\groupmail}{alittlebyte.swe@gmail.com}
\newcommand{\grouplogo}{Immagini/logo_alittlebyte.png}
\newcommand{\groupsymbol}{Immagini/solo_logo.png}

\geometry{
    left=2.5cm,
    right=2.5cm,
    top=2.5cm,
    bottom=2.5cm,
    headheight=331pt,
    includeheadfoot
}

\pagestyle{fancy}
\fancyhf{}

\renewcommand{\headrulewidth}{0pt}
\fancyhead{}

\renewcommand{\footrulewidth}{0.4pt}
\fancyfoot[L]{\includegraphics[height=0.6cm]{\groupsymbol}\hspace{0.45em}aLittleByte}
    \fancyfoot[C]{\thepage}
\fancyfoot[R]{Verbale 2026-04-27}

\fancypagestyle{firstpage}{
    \fancyhf{}
    \renewcommand{\headrulewidth}{0pt}
    \renewcommand{\footrulewidth}{0.4pt}
    \fancyhead[C]{%
        \makebox[\textwidth][c]{%
            \begin{tabular}{c}
                \includegraphics[height=10cm]{\grouplogo}\\[1ex]
                \small\texttt{\groupmail}
            \end{tabular}%
        }
    }
    \fancyfoot[L]{\includegraphics[height=0.6cm]{\groupsymbol}\hspace{0.45em}aLittleByte}
        \fancyfoot[C]{\thepage}
    \fancyfoot[R]{Verbale 2026-04-27}
}

\begin{document}

\thispagestyle{firstpage}

\vspace*{0.5cm}

\begin{center}
    {\LARGE \textbf{Verbale Interno - 2026-04-27}}
\end{center}

\vspace{1.5cm}

\begin{center}
    \renewcommand{\arraystretch}{1.3}
    \begin{tcolorbox}[
        width=0.92\textwidth,
        colback=gray!6,
        colframe=gray!45,
        boxrule=0.5pt,
        arc=2mm,
        boxsep=0pt,
        left=8pt, right=8pt, top=8pt, bottom=8pt
    ]
        \begin{tabularx}{\linewidth}{>{\bfseries}p{3.5cm} X}
            Gruppo      & aLittleByte (gruppo 19) \\
            Redattore   & Eleonora Bellet \\
            Verificatori & Leonardo Soligo\\
            Destinatari & Prof. Tullio Vardanega, Prof. Riccardo Cardin \\
            Data        & 27 Aprile 2026\\
        \end{tabularx}
    \end{tcolorbox}
\end{center}

\newpage

\newgeometry{left=2.5cm,right=2.5cm,top=2.5cm,bottom=2.5cm,includefoot}
\tableofcontents
\newpage
\section{Informazioni Generali}
\section*{Specifiche Documento}
\hrule
\vspace{1cm}

\begin{tabular}{>{\bfseries}p{5cm} | p{10cm}}
Luogo Riunione & Discord \\[6pt]

Orario (Inizio - Fine) & 15:30 -- 16.15 \\[6pt]

Redattore &
  E. Bellet\\[6pt]

Amministratore & 
A. Gastaldello\\[6pt]

Verificatore &
L. Soligo\\[6pt]

Partecipanti &
  A. Oloni\\ &
  L. Soligo\\&
  E. Bellet\\&
  A. Gastaldello\\&
  A. Maule\\&
  V. Y. Kaneda Fernandes\\&
  D. Belluz\\

\end{tabular}


\section{Ordine del Giorno}
    \begin{itemize}
    \item Stato di avanzamento dei casi d'uso e dei documenti di progetto
    \item Definizione del modello di sviluppo e organizzazione degli sprint
    \item Gestione e conteggio delle ore nei vari sprint
    \item Organizzazione e aggiornamento del Glossario
    \item Continuazione e gestione del Mock-up
    \end{itemize}


\section{Attività svolte}
\subsection{Stato di avanzamento dei casi d'uso e dei documenti di progetto}
Il team ha effettuato un aggiornamento generale sui documenti di progetto, confermando il completamento dei casi d'uso per l'Analisi dei Requisiti, alla quale mancano solamente i diagrammi. È stata confermata anche la sistemazione delle Norme di Progetto. La sezione relativa alla qualità del software e ai test è stata rimandata a una fase successiva.

\subsection{Definizione del modello di sviluppo e organizzazione degli sprint}
È stata discussa la scelta del ciclo di vita del software, confrontando il modello incrementale con quello Agile. La decisione finale è ricaduta sull'adozione del modello Agile, basato su sprint bisettimanali, come già scelto in precedenza. Si è inoltre stabilito di tenere traccia puntuale delle date e dell'organizzazione degli sprint all'interno del piano di progetto.

\subsection{Gestione e conteggio delle ore nei vari sprint}
Il team ha esaminato la distribuzione delle ore lavorate. Per rendere il conteggio più flessibile rispetto alla modifica diretta sul piano di progetto, si è deciso di utilizzare un foglio Excel condiviso. È stata ribadita l'importanza di allineare le ore al piano di qualifica dei ruoli, con l'obiettivo di far raggiungere a ciascun membro un totale di 90 ore lavorative entro il mese di settembre. Il responsabile del piano di progetto si farà carico di questa gestione su Excel.

\subsection{Organizzazione e aggiornamento del Glossario}
È stata definita una procedura condivisa per l'aggiornamento del glossario: i termini tecnici o ambigui verranno identificati direttamente durante la revisione dei documenti. Inoltre, il team sta valutando l'introduzione di automatismi, tramite strumenti di LaTeX o GitHub, per collegare automaticamente le parole presenti nei documenti al glossario.

\subsection{Continuazione e gestione del Mock-up}
Si è discusso del Mock-up realizzato da \textbf{Leonardo Soligo}, stabilendo che necessita di una verifica accurata per assicurare l'allineamento con le indicazioni dell'analisi dei requisiti.

\newpage
\section{To-Do List}
\begin{table}[h]
    \centering
    \begin{tabular}{|p{4.5cm}|p{5cm}|l|}
    \hline
        \rowcolor{lightgray}\textit{\textbf{Persona Incaricata}}  & \textbf{Task Assegnata} &\textbf{Link Issue} \\
    \hline
        \textit {Leonardo Soligo} &  Presentazione del Mock-up per l'azienda & \href{https://github.com/aLittleByte-19/Documentazione/issues/74}{\#74}\\
    \hline
        \textit{Angelica Gastaldello e Angela Maule} &  Completare i diagrammi per l'Analisi dei Requisiti& \href{https://github.com/aLittleByte-19/Documentazione/issues/51}{\#51}\\
    \hline
        \textit{Vinicius Yuri Kaneda Fernandes e Diego Belluz} & Continuazione Stesura Glossario & \href{https://github.com/aLittleByte-19/Documentazione/issues/53}{\#53}\\
    \hline
        \textit{Alex Oloni e Eleonora Bellet} & Continuazione Stesura Piano di Progetto &
        \href{https://github.com/aLittleByte-19/Documentazione/issues/63}{\#63}\\
    \hline
        \textit{Alex Oloni} & Creare il file Excel condiviso e aggiornare le ore di lavoro & \href{https://github.com/aLittleByte-19/Documentazione/issues/70}{\#70}\\
    \hline
        \textit{Angelica Gastaldello} & Scrittura email alla azienda per la prossima riunione & 
        \href{https://github.com/aLittleByte-19/Documentazione/issues/71}{\#71}\\
    \hline 
        \textit{Tutti i pertecipanti} & Domande da porre alla azienda per la prossima riunione& 
        \href{https://github.com/aLittleByte-19/Documentazione/issues/72}{\#72}\\
    \hline
        \textit{Eleonora Bellet} & Stesura Verbale Interno 27/04/2026 & \href{https://github.com/aLittleByte-19/Documentazione/issues/73}{\#73}\\
    \hline
    \end{tabular}
    \label{tab:placeholder}
\end{table}

\end{document}